% 文献目录 Bibliography (原书页 183–188 / PDF p198–p203)
% PDF p204 为封底（仅 ISBN 981-02-3231-4(pbk) 条形码），无正文文字。
% 编号注：原书在 6 与 7 之间有一条 "6a."（M. Lax），故 \item 共 108 行，
% 但编号与原书完全一致：(1)…(6), (6a), (7)…(107)。6a 条目用 \item[(6a)]
% 手工覆盖标签；enumitem 对显式标签不再步进计数器，故后续 \item 自动
% 从 (7) 起、与原书编号对齐，无需手动调整计数器。
% 照录原书排印：Frohlich（无变音符）、Hulthen（无重音）、Le Problem a N Corps
% （法文书名原书即无重音）、J.E.T.P.（67/73 条紧凑）与 J. E. T. P.（88 条带空格）
% 两种写法并存、期刊缩写句点均照录。
\chapter*{文献目录}
\addcontentsline{toc}{chapter}{文献目录}

\begin{enumerate}[label=(\arabic*),leftmargin=2.8em,itemsep=1pt,parsep=0pt]
\item C. Kittel, ``Introduction to Solid State Physics,'' 2nd ed., Wiley, New York, 1953.
\item F. Seitz, ``Modern Theory of Solids,'' McGraw-Hill, New York, 1940.
\item J. M. Ziman, ``Electrons and Phonons,'' Clarendon Press, Oxford, 1960.
\item G. H. Wannier, ``Elements of Solid State Theory,'' Cambridge University Press, Cambridge, 1959.
\item R. Peierls, ``Quantum Theory of Solids,'' Clarendon Press, Oxford, 1955.
\item L. D. Landau and E. M. Lifshitz, ``Statistical Physics,'' Pergamon, London, 1958.
\item[(6a)] M. Lax, ``Symmetry Principles in Solid State Physics,'' to be published (publisher undecided).
\item L. D. Landau, private communication.
\item This discussion is from A. M. Clogston and V. Jaccarino, \emph{Phys. Rev.}, 121, 1357 (1961).
\item F. J. Morin, \emph{Bell System Tech. J.}, 37, 1047 (1958); \emph{Phys. Rev. Letters}, 3, 34 (1959).
\item A good review is: V. B. Compton, T. H. Geballe, and B. T. Matthias, \emph{Rev. Mod. Phys.}, to be published.
\item John R. Reitz, ``Methods of the One-Electron Theory of Solids,'' in Solid State Physics, Advances in Research and Applications, F. Seitz and D. Turnbull (eds.), Academic, New York, 1955, Vol. 1.
\item H. Ehrenreich and M. H. Cohen, \emph{Phys. Rev.}, 115, 786 (1959); for preceding work see references therein.
\item D. J. Thouless and J. G. Valatin, \emph{Phys. Rev. Letters}, 5, 509 (1960).
\item For instance, J. G. Valatin, \emph{Phys. Rev.}, 122, 1012 (1961).
\item G. W. Pratt, Jr., \emph{Phys. Rev.}, 102, 1303 (1956); W. Marshall, \emph{Proc. Phys. Soc. (London)}, 78, 113 (1961).
\item J. C. Slater, \emph{J. Chem. Phys.}, 19, 220 (1951).
\item D. J. Thouless, ``The Quantum Mechanics of Many-Body Systems,'' pp. 24-29, Academic, London, 1961.
\item The first reference I know to this as a general quantum mechanical procedure is F. J. Dyson, \emph{Phys. Rev.}, 90, 994; 91, 1543 (1953). He calls it the ``new Tamm-Dancoff'' method.
\item F. S. Ham, ``The Quantum Defect Method,'' in Solid State Physics, Advances in Research and Applications, F. Seitz and D. Turnbull (eds.), Academic, New York, 1955, Vol. 1.
\item E. P. Wigner and F. Seitz, ``Qualitative Analysis of the Cohesion of Metals,'' in Solid State Physics, Advances in Research and Applications, F. Seitz and D. Turnbull (eds.), Academic, New York, 1955, Vol. 1.
\item H. Brooks, \emph{Nuovo Cimento Suppl.}, [X] 7, 165 (1958).
\item H. Jones, ``Theory of Brillouin Zones and Electronic States in Crystals,'' North-Holland, Amsterdam, 1960.
\item The first to consider that nearly free electron theory might be relevant to the diamond lattices was F. Herman, \emph{Phys. Rev.}, 93, 1214 (1954).
\item A. W. Overhauser, \emph{Phys. Rev.}, 128, 1437 (1962).
\item J. C. Phillips and L. Kleinman, \emph{Phys. Rev.}, 116, 880 (1959); 118, 1153 (1960).
\item H. Jones (Ref. 22), Chap. 5, Sec. 44.
\item J. Bardeen, \emph{J. Chem. Phys.}, 6, 367 (1938).
\item E. P. Wigner and F. Seitz, \emph{Phys. Rev.}, 46, 509 (1934).
\item E. P. Wigner, \emph{Phys. Rev.}, 46, 1002; \emph{Trans. Faraday Soc.}, 34, 678.
\item P. A. M. Dirac, \emph{Proc. Cambridge Phil. Soc.}, 26, 376 (1930).
\item Probably the best source is D. Pines, in Solid State Physics, Advances in Research and Applications, F. Seitz and D. Turnbull (eds.), Academic, New York, 1955, Vol. 1, p. 367.
\item J. C. Slater, \emph{Phys. Rev.}, 81, 385 (1951).
\item L. Kleinman and J. C. Phillips, \emph{Phys. Rev.}, 116, 380 (1959). This is an interesting numerical illustration of many of the ideas of the next chapter. Also note the discussion of exchange in the appendix.
\item C. Herring, \emph{Phys. Rev.}, 57, 1169 (1940).
\item C. Herring and A. G. Hill, \emph{Phys. Rev.}, 58, 132 (1940).
\item M. L. Cohen and V. Heine, \emph{Phys. Rev.}, 122, 1821 (1961).
\item B. J. Austin, V. Heine, and L. J. Sham, \emph{Phys. Rev.}, 127, 276 (1962).
\item See, for instance, K. A. Brueckner in ``The Many-Body Problem,'' Dunod-Wiley, Paris, New York, 1959.
\item For example: J. M. Ziman, \emph{Phil. Mag.}, 6, 1013 (1961); C. C. Bradley, T. E. Faber, E. G. Wilson, and J. M. Ziman, \emph{Phil. Mag.}, 7, 865 (1962).
\item B. T. Matthias, \emph{Phys. Rev.}, 97, 74 (1955).
\item C. Herring, \emph{J. Appl. Phys.}, \emph{Suppl.}, 31, 3S (1960).
\item H. A. Bethe and E. Salpeter, ``Quantum Mechanics of One- and Two-Electron Systems,'' Academic, New York, 1957, p. 36.
\item J. C. Slater, \emph{Phys. Rev.}, 36, 57 (1930).
\item W. Kohn, ``Shallow Impurity States in Si and Ge,'' in Solid State Physics, Advances in Research and Applications, F. Seitz and D. Turnbull (eds.), Academic, New York, 1957, Vol. 5, p. 258.
\item E. I. Blount, ``Formalisms of Band Theory,'' in Solid State Physics, Advances in Research and Applications, F. Seitz and D. Turnbull (eds.), Academic, New York, 1962, Vol. 13, p. 305.
\item Throughout this chapter I shall omit references for the various effects unless they are of particular interest in the exposition, or not well covered in the various texts. Explanations and original references will be found in the texts and reviews referred to. Another good source for free-electron effects is A. B. Pippard, ``Experimental Analysis of the Electronic Structure of Metals,'' \emph{Rept. Progr. Phys.}, 23, 176 (1960).
\item G. Feher, \emph{Phys. Rev.}, 114, 1219 (1959). Note Fig. 8; as far as I know, this is the only published version of the theoretical results on the interference effect.
\item G. H. Wannier (Ref. 4), Chap. 6.
\item A. G. Chynoweth, G. H. Wannier, R. A. Logan, and D. E. Thomas, \emph{Phys. Rev. Letters}, 5, 57 (1960).
\item See especially Pippard, Ref. 46.
\item A good review of the open-orbit situation is given by R. G. Chambers in his article on ``Magnetoresistance'' in ``The Fermi Surface,'' W. A. Harrison and M. B. Webb (eds.), Wiley, New York, 1960, p. 100. He refers to the important Russian work.
\item J. M. Luttinger, \emph{Phys. Rev.}, 102, 1030 (1956).
\item R. C. Fletcher, W. A. Yager, and F. R. Merritt, \emph{Phys. Rev.}, 100, 747 (1955).
\item M. H. Cohen and L. Falicov, \emph{Phys. Rev. Letters}, 7, 231 (1961).
\item E. I. Blount, \emph{Phys. Rev.}, 126, 1636 (1962). Another excellent paper giving a more complete theory is A. B. Pippard, \emph{Proc. Roy. Soc. (London)}, A270, 1 (1962).
\item References and a more complete discussion will be found in Blount's review (Ref. 45).
\item R. Karplus and J. M. Luttinger, \emph{Phys. Rev.}, 95, 1154 (1954); J. M. Luttinger, \emph{Phys. Rev.}, 112, 739 (1958).
\item D. Pines, ``The Many-Body Problem,'' Benjamin, New York, 1961.
\item D. J. Thouless, ``Quantum Mechanics of Many-Body Systems,'' Academic, London, 1961.
\item C. DeWitt (ed.), ``The Many-Body Problem,'' Methuen, London, and Wiley, New York, 1959.
\item V. F. Weisskopf, \emph{Science}, 113, 101 (1951).
\item R. MacWeeny, \emph{Proc. Roy. Soc. (London)}, A223, 63, 306 (1954).
\item L. D. Landau, \emph{J. Phys. U.S.S.R.}, 5, 71 (1941).
\item Frohlich, Pelzer, and Zienau, \emph{Phil. Mag.}, 41, 221 (1950).
\item T. D. Lee and D. Pines, \emph{Phys. Rev.}, 88, 960 (1952).
\item W. Kohn, \emph{Phys. Rev.}, 105, 509 (1957).
\item J. M. Luttinger, \emph{Phys. Rev.}, 119, 1153 (1960). See also A. B. Migdal, \emph{J.E.T.P.}, 32, 399 (1957); \emph{Soviet Phys. (J.E.T.P.)}, 5, 333 (1957).
\item The development here is close to that of P. Nozi\`eres' lecture notes, ``Le Problem a N Corps,'' Dunod, Paris, 1963 (trans.: ``Theory of Interacting Fermi Systems,'' Benjamin, New York, 1963). It is also used in P. Nozi\`eres and J. M. Luttinger, \emph{Phys. Rev.}, 127, 1431 (1962).
\item W. Kohn, \emph{Phys. Rev.}, 110, 857 (1958); see also V. Ambegaokar, \emph{Phys. Rev.}, 121, 91 (1961).
\item J. Bardeen and W. Shockley, \emph{Phys. Rev.}, 80, 72 (1950); for generalization see C. Herring and W. Vogt, \emph{Phys. Rev.}, 101, 944 (1956); Ziman's book; and Blount's article.
\item This situation was first discussed physically by R. R. Heikes and W. D. Johnston, \emph{J. Chem. Phys.}, 26, 582 (1957); a more complete theoretical treatment is given by, e.g., T. Holstein, \emph{Ann. Phys.}, 8, 325, 343 (1959).
\item H. Frohlich, \emph{Phys. Rev.}, 79, 845 (1950); J. Bardeen, \emph{Phys. Rev.}, 80, 567 (1950).
\item L. D. Landau, \emph{J.E.T.P.}, 30, 1058 (1956) [trans.: \emph{Soviet Phys. (J.E.T.P.)}, 3, 920]. An excellent review of Fermi liquid theory which follows Landau's point of view is A. A. Abrikosov and I. M. Khalatnikov, \emph{Usp. Fiz. Nauk}, 56, 177 (1958) [trans.: \emph{Soviet Phys. Usp. (London)}, 66, 68 (1958)].
\item J. Goldstone, \emph{Proc. Roy. Soc. (London)}, A293, 267 (1957); W. Kohn and J. M. Luttinger, \emph{Phys. Rev.}, 118, 41 (1960); J. M. Luttinger and J. C. Ward, \emph{Phys. Rev.}, 118, 1417 (1960); J. M. Luttinger, \emph{Phys. Rev.}, 119, 1153 (1960).
\item J. M. Luttinger and P. Nozi\`eres, \emph{Phys. Rev.}, 127, 1425, 1431 (1962).
\item This difficulty is mentioned by Abrikosov and Khalatnikov (Ref. 73), App. 1. Most striking is the failure associated with the B.C.S. theory of superconductivity.
\item J. Frenkel, \emph{Phys. Rev.}, 17, 17 (1931); \emph{Phys. Z. Sowjetunion}, 8, 185 (1935).
\item The Wannier exciton seems basically to have been introduced in an excellent paper by J. C. Slater and W. Shockley, \emph{Phys. Rev.}, 50, 705 (1936), although Wannier's paper [G. H. Wannier, \emph{Phys. Rev.}, 52, 191 (1937)] introduces the idea of Wannier functions and the effective hamiltonian as a formal theoretical scheme.
\item J. J. Hopfield, \emph{Phys. Rev.}, 112, 1555 (1958).
\item E. Dresselhaus, \emph{J. Phys. Chem. Solids}, 1, 14 (1956).
\item F. J. Dyson, \emph{Phys. Rev.}, 102, 1217, 1230 (1956).
\item G. Wentzel, \emph{Phys. Rev.}, 108, 1593 (1957).
\item See also T. Matsubara and J. M. Blatt, \emph{Progr. Theoret. Phys.}, 23, 451 (1960), and other papers by the same authors for another approach.
\item M. H. Cohen and F. Keffer, \emph{Phys. Rev.}, 99, 1128 (1955). For the exciton application see W. R. Heller and A. Marcus, \emph{Phys. Rev.}, 84, 809 (1951).
\item R. H. Lyddane, R. G. Sachs, and E. Teller, \emph{Phys. Rev.}, 59, 673 (1941).
\item D. G. Thomas and J. J. Hopfield, \emph{Phys. Rev.}, 124, 657 (1961).
\item N. N. Bogolyubov, \emph{J. Phys. U.S.S.R.}, 9, 23 (1947).
\item For methods similar to these see N. N. Bogolyubov, \emph{J. E. T. P.}, 34, 73 (1958); and Thouless' book.
\item A discussion of this, referring also to Ferrell and others, is given by G. E. Brown and D. J. Thouless, \emph{Physica, Suppl.}, 26, S145 (1960).
\item For these ideas it is better to read one of the books on many-body theory referred to [ (58-60); also \emph{Physica, Suppl.}, 26, has a number of good semi-review papers ] than to attempt to study the confusing and profuse original literature. The names Brueckner, Watson, and Eden are especially associated with ``ladder'' summing; Pines, Nozi\`eres, Hubbard, and Ferrell with ``ring'' sums.
\item For a review giving an expanded version of the viewpoint used here, see P. W. Anderson, in Solid State Physics, Advances in Research and Applications, F. Seitz and D. Turnbull (eds.), Academic, New York, 1963, Vol. 14, p. 99; also most of the relevant references may be found there.
\item J. C. Slater, \emph{Phys. Rev.}, 52, 198 (1937).
\item S. Schubin and S. Vonsovsky, \emph{Proc. Roy. Soc. (London)}, 145, 159 (1934).
\item J. C. Slater, \emph{Rev. Mod. Phys.}, 25, 199 (1953).
\item C. Herring, \emph{Rev. Mod. Phys.}, 34, 631 (1962).
\item A brief review of spin waves will be found in Ref. 91. More complete is a forthcoming article by F. Keffer in ``Magnetism,'' G. T. Rado and H. Suhl (eds.), Addison-Wesley, Reading, Mass., in press.
\item C. J. Gorter and J. H. Van Vleck, \emph{Phys. Rev.}, 72, 1128 (1947).
\item Considerations of this sort are discussed by Landau and Lifshitz (Ref. 6), Chap. 13, Sec. 125.
\item M. J. Klein and R. S. Smith, \emph{Phys. Rev.}, 80, 1111 (1951).
\item C. Herring and C. Kittel, \emph{Phys. Rev.}, 81, 869 (1951); L. D. Landau and E. M. Lifshitz, \emph{Phys. Z. Sowjetunion}, 8, 153 (1935).
\item P. W. Anderson, \emph{Phys. Rev.}, 86, 694 (1952); L. Hulthen, \emph{Proc. Roy. Acad. Sci. Amst.}, 39, 190 (1936). A review is T. Nagamiya, R. Kubo, and K. Yosida, \emph{Phil. Mag., Suppl.}, 14, 1 (1955).
\item A name introduced by J. Goldstone, A. Salam, and S. Weinberg, \emph{Phys. Rev.}, 127, 965 (1962).
\item L. Hulthen, \emph{Ark. Mat. Astron. Fys.}, 26A, No. 1 (1938).
\item P. W. Anderson and D. J. Thouless, \emph{Phys. Letters}, 1, 155 (1962).
\item N. N. Bogolyubov, \emph{Physica, Suppl.}, 26, S1 (1960).
\item L. R. Walker, \emph{Phys. Rev.}, 105, 390 (1957).
\item P. W. Anderson, \emph{Phys. Rev.}, 112, 1900 (1958).
\end{enumerate}
